\documentclass[11pt]{article}
% ============ FangMath self-contained preamble (pdflatex) ============
\usepackage[T1]{fontenc}
\usepackage{mathpazo}            % Palatino text + math
\usepackage[margin=1in]{geometry}
\usepackage{amsmath,amssymb}
\usepackage{xcolor}
\usepackage[most]{tcolorbox}
\usepackage{fancyhdr}
\usepackage{enumitem}
\usepackage{pifont}
\usepackage{needspace}
\usepackage{tikz}
\usetikzlibrary{arrows.meta}
\definecolor{FangBlue}{HTML}{1F4E79}
\definecolor{FangLight}{HTML}{EAF1F8}
\definecolor{FangGreen}{HTML}{1D9E75}
\definecolor{FangRed}{HTML}{C0392B}
% ---- Headers / footers ----
\setlength{\headheight}{26pt}
\pagestyle{fancy}
\fancyhf{}
\renewcommand{\headrulewidth}{0.8pt}
\renewcommand{\headrule}{\hbox to\headwidth{\color{FangBlue}\leaders\hrule height \headrulewidth\hfill}}
\lhead{\color{FangBlue}\textbf{FangMath}}
\rhead{\color{FangBlue}Absolute Value Inequalities}
\cfoot{\color{FangBlue}\thepage}
% ---- Step box ----
\newtcolorbox{stepbox}[1]{
  colback=FangLight, colframe=FangBlue, boxrule=0.8pt, arc=3pt,
  left=10pt, right=10pt, top=6pt, bottom=6pt,
  fonttitle=\bfseries\color{white},
  colbacktitle=FangBlue, coltitle=white,
  title={#1}, breakable}
% ---- Mistake box ----
\newtcolorbox{mistakebox}{
  colback=white, colframe=FangRed, boxrule=1pt, arc=3pt,
  left=10pt, right=10pt, top=6pt, bottom=6pt,
  fonttitle=\bfseries\color{white},
  colbacktitle=FangRed, coltitle=white,
  title={Common Mistake to Avoid}, breakable}
% ---- Answer box ----
\newtcolorbox{answerbox}{
  colback=white, colframe=FangGreen, boxrule=1.2pt, arc=3pt,
  left=10pt, right=10pt, top=6pt, bottom=6pt}
\newcommand{\In}{\textbf{\color{FangGreen}in}}
\newcommand{\Out}{\textbf{\color{FangRed}out}}
% ---- Number line: just the two boundary dots, no shading yet ----
\newcommand{\numlineDots}[6]{%
\begin{center}
\begin{tikzpicture}[x=0.42cm]
  \draw[thick,-{Stealth[length=2mm]}] (#1-0.5,0) -- (#2+0.5,0);
  \draw[thick,-{Stealth[length=2mm]}] (#2+0.5,0) -- (#1-0.5,0);
  \filldraw[FangBlue] (#3,0) circle (3.2pt);
  \filldraw[FangBlue] (#4,0) circle (3.2pt);
  \node[below=5pt] at (#3,0) {#5};
  \node[below=5pt] at (#4,0) {#6};
\end{tikzpicture}
\end{center}}
% ---- Number line: shaded interval BETWEEN two closed endpoints (And / In) ----
% #1=min #2=max #3=left endpoint #4=right endpoint #5=left label #6=right label
\newcommand{\numlineIn}[6]{%
\begin{center}
\begin{tikzpicture}[x=0.42cm]
  \draw[thick,-{Stealth[length=2mm]}] (#1-0.5,0) -- (#2+0.5,0);
  \draw[thick,-{Stealth[length=2mm]}] (#2+0.5,0) -- (#1-0.5,0);
  \draw[line width=3pt,FangGreen] (#3,0) -- (#4,0);
  \filldraw[FangGreen] (#3,0) circle (3.2pt);
  \filldraw[FangGreen] (#4,0) circle (3.2pt);
  \node[below=5pt] at (#3,0) {#5};
  \node[below=5pt] at (#4,0) {#6};
\end{tikzpicture}
\end{center}}
% ---- Number line: shaded rays OUTSIDE two closed endpoints (Or / Out) ----
\newcommand{\numlineOut}[6]{%
\begin{center}
\begin{tikzpicture}[x=0.42cm]
  \draw[thick,-{Stealth[length=2mm]}] (#1-0.8,0) -- (#2+0.8,0);
  \draw[thick,-{Stealth[length=2mm]}] (#2+0.8,0) -- (#1-0.8,0);
  \draw[line width=3pt,FangGreen] (#1-0.8,0) -- (#3,0);
  \draw[line width=3pt,FangGreen] (#4,0) -- (#2+0.8,0);
  \filldraw[FangGreen] (#3,0) circle (3.2pt);
  \filldraw[FangGreen] (#4,0) circle (3.2pt);
  \node[below=5pt] at (#3,0) {#5};
  \node[below=5pt] at (#4,0) {#6};
\end{tikzpicture}
\end{center}}
\begin{document}
\begin{center}
  {\color{FangBlue}\LARGE\bfseries Solving Absolute Value Inequalities}\\[3pt]
  {\color{FangBlue}\large The ``In or Out'' Method}
\end{center}
\vspace{4pt}

\begin{stepbox}{The Big Rule}
Absolute value bars ask: how far is the inside expression from $0$? Instead of memorizing which symbol
means \emph{and} and which means \emph{or}, \textbf{test the middle point} (where the inside expression
equals $0$) to see whether it belongs to the solution:
\begin{center}
\renewcommand{\arraystretch}{1.3}
\begin{tabular}{|c|c|c|}
\hline
\textbf{Test result at the middle} & \textbf{Region} & \textbf{Compound inequality} \\\hline
True & \In\ (between the boundaries) & \textbf{and} \\\hline
False & \Out\ (outside the boundaries) & \textbf{or} \\\hline
\end{tabular}
\end{center}
This works because $\lvert \text{expr}\rvert \le k$ describes a single connected interval around $0$
(an \emph{and}), while $\lvert \text{expr}\rvert \ge k$ describes everything \emph{except} that interval
(an \emph{or}).
\end{stepbox}

\vspace{4pt}
\noindent{\color{FangBlue}\Large\bfseries Example 1 \quad Solve $\lvert 3x-4\rvert \le 8$ \normalsize (an ``And'' case)}
\vspace{2pt}

\begin{stepbox}{Step 1 \ Find the boundaries}
Since $\lvert -8\rvert = 8$ and $\lvert 8\rvert = 8$, the expression $3x-4$ hits the edge at
$-8$ and $8$. Mark both as \textbf{closed} dots (the inequality is $\le$, so the ends are included).
\numlineDots{-8}{8}{-8}{8}{$-8$}{$8$}
\end{stepbox}
\begin{stepbox}{Step 2 \ Check: In or Out?}
Test the middle point $3x-4 = 0$:
\[
\lvert 0\rvert = 0 \le 8 \quad\text{\textbf{\color{FangGreen}True}.}
\]
The point \emph{between} the boundaries works, so we keep the \In\ region. That makes this an
\textbf{and} statement (a \emph{between} inequality), not an \emph{or}.
\end{stepbox}
\begin{stepbox}{Step 3 \ Write the compound inequality and solve}
Because it is ``in,'' $3x-4$ stays \emph{between} $-8$ and $8$:
\[
-8 \le 3x-4 \le 8.
\]
Solve all three parts at once. Add $4$:
\[
-8+4 \le 3x-4+4 \le 8+4
\qquad\Longrightarrow\qquad
-4 \le 3x \le 12.
\]
Divide by $3$:
\[
\frac{-4}{3} \le x \le \frac{12}{3}
\qquad\Longrightarrow\qquad
-\frac{4}{3} \le x \le 4.
\]
\end{stepbox}
\begin{center}
\begin{answerbox}
\centering
\textbf{Solution:} \quad $\displaystyle -\frac{4}{3} \le x \le 4$
\quad or \quad $\left[-\dfrac{4}{3},\, 4\right]$
\end{answerbox}
\end{center}
\numlineIn{-2}{5}{-1.333}{4}{$-\frac{4}{3}$}{$4$}
\vspace{2pt}
\begin{center}
\small\itshape Check: at $x=0$, $\lvert 3(0)-4\rvert = 4 \le 8$ \checkmark
\end{center}

\vspace{10pt}
\noindent{\color{FangBlue}\Large\bfseries Example 2 \quad Solve $\lvert 2x+4\rvert \ge 6$ \normalsize (an ``Or'' case)}
\vspace{2pt}

\noindent Same method as Example 1 -- but this time the test point will fail, which flips it to an
\emph{or} statement.

\begin{stepbox}{Step 1 \ Find the boundaries}
Since $\lvert -6\rvert = 6$ and $\lvert 6\rvert = 6$, the expression $2x+4$ hits the edge at
$-6$ and $6$. Mark both as \textbf{closed} dots (the inequality is $\ge$, so the ends are included).
\numlineDots{-6}{6}{-6}{6}{$-6$}{$6$}
\end{stepbox}
\begin{stepbox}{Step 2 \ Check: In or Out?}
Test the middle point $2x+4 = 0$:
\[
\lvert 0\rvert = 0 \ge 6 \quad\text{\textbf{\color{FangRed}False}.}
\]
The point \emph{between} the boundaries does \textbf{not} work, so the solution is the \Out\ region.
That makes this an \textbf{or} statement (an \emph{outside} inequality), not an \emph{and}.
\end{stepbox}
\begin{stepbox}{Step 3 \ Write the compound inequality and solve}
Because it is ``out,'' $2x+4$ stays \emph{outside} $-6$ and $6$ --- less than or equal to the left
boundary, \textbf{or} greater than or equal to the right boundary:
\[
2x+4 \le -6 \qquad\text{\textbf{or}}\qquad 2x+4 \ge 6.
\]
Solve each piece on its own. Subtract $4$ from both sides:
\[
2x \le -10 \qquad\text{\textbf{or}}\qquad 2x \ge 2.
\]
Divide both sides by $2$:
\[
x \le -5 \qquad\text{\textbf{or}}\qquad x \ge 1.
\]
\end{stepbox}

\needspace{9\baselineskip}
\begin{mistakebox}
The original inequality is $\ge$, not a strict $>$ -- so the boundary points themselves are part of the
solution, and the compound inequality must use $\le$ and $\ge$, not $<$ and $>$. Writing ``$2x+4<-6$ or
$2x+4>6$'' would quietly kick $x=-5$ and $x=1$ back out of the answer. Always check: plug the boundary
value into the \emph{original} inequality --- at $x=-5$, $\lvert 2(-5)+4\rvert = \lvert -6\rvert = 6 \ge 6$
\checkmark, so it belongs, and the dots on the number line stay \textbf{closed}.
\end{mistakebox}

\begin{center}
\begin{answerbox}
\centering
\textbf{Solution:} \quad $x \le -5 \ \text{ or } \ x \ge 1$
\quad or \quad $(-\infty,-5\,] \cup [\,1,\infty)$
\end{answerbox}
\end{center}
\numlineOut{-9}{4}{-5}{1}{$-5$}{$1$}
\vspace{2pt}
\begin{center}
\small\itshape Check: at $x=-10$ (outside), $\lvert 2(-10)+4\rvert = 16 \ge 6$ \checkmark \qquad
at $x=0$ (between), $\lvert 2(0)+4\rvert = 4 \ge 6$? \ding{55}\ (correctly excluded)
\end{center}

\vspace{6pt}
\begin{stepbox}{Key Idea --- And vs.\ Or at a Glance}
\begin{center}
\renewcommand{\arraystretch}{1.3}
\begin{tabular}{|c|c|c|}
\hline
\textbf{Form} & \textbf{Test at the middle} & \textbf{Solution shape}\\\hline
$\lvert \text{expr}\rvert \le k$ & always True (In) & \textbf{and}: one connected interval\\\hline
$\lvert \text{expr}\rvert \ge k$ & always False (Out) & \textbf{or}: two separate rays\\\hline
\end{tabular}
\end{center}
(Assuming $k>0$.) The middle-point test is a shortcut for the graph of $y=\lvert\text{expr}\rvert$: it
dips \emph{below} the line $y=k$ only in the middle (giving an \emph{and}), so anywhere it stays at or
above $y=k$ -- like $\lvert 2x+4\rvert$ compared to $y=6$ in Example 2 -- gives an \emph{or}.
\end{stepbox}

\end{document}
